\documentclass{article}

\textwidth=6.5in
\textheight=8.5in
\voffset=-54pt
\oddsidemargin=0in
\evensidemargin=0in
\setlength{\parskip}{8pt}
\setlength{\parindent}{0pt}

\usepackage{amsmath, amssymb}
\usepackage{ifthen}

\usepackage[inline]{enumitem}
\setlist{topsep=0pt, parsep=8pt}

\usepackage[rgb,table]{xcolor}\convertcolorsUfalse
\usepackage{graphicx}

% change hyperref options
\PassOptionsToPackage{hyphens}{url}
\usepackage[bookmarks,colorlinks,linkcolor=black,citecolor=blue,pdfstartview=FitH,urlcolor=blue]{hyperref}
\hypersetup{pdfpagemode=UseNone}
\usepackage{bookmark}

\usepackage{graphicx}
\usepackage{multicol}
\usepackage{framed}
\usepackage{array}

\usepackage{icomma}

\usepackage{pifont}

\usepackage{marginnote}
\usepackage[colorinlistoftodos]{todonotes}
\renewcommand{\marginpar}{\marginnote}

\usepackage{soul}

\usepackage{pgfplots}
\pgfplotsset{compat=1.18}  
\pgfplotscreateplotcyclelist{mycolor}{black, blue, red, green!70!black, magenta, purple, cyan, orange }  
\pgfplotsset{
  disabledatascaling,
  cycle list=tikzcolor, 
  axis lines=middle,
  every axis plot/.style={no marks, thick},
  every axis/.style={
    enlarge x limits=false,
    enlarge y limits=auto,
    xlabel style={at={(current axis.right of origin)},anchor=west},
    ylabel style={at={(current axis.above origin)},anchor=south},
    % extra x and y ticks for asymptotes
    extra x tick style={grid=major, grid style={dashed,black}},
    extra y tick style={grid=major, grid style={dashed,black}},
    /pgf/number format/1000 sep={},
  },    
  tick label style = {font=\footnotesize, fill=\currentbgcolor, inner sep=0pt},    
  xticklabel shift=1pt,    
  scaled ticks=false,
  scaled y ticks=false,
  unbounded coords=jump,
  samples=101,
  minor tick length=0,
  scatter/classes={
    open={mark=*,draw=black,fill=white, mark size=2, thin},
    closed={mark=*,draw=black,fill=black, mark size=2, thin}
  },
  width=3in,
}
\pgfplotsset{open/.style={only marks, mark=*, fill=white, thin, mark size=2}}
\pgfplotsset{closed/.style={only marks, mark=*, draw=black, fill=black,
    thin, mark size=2}}
\pgfplotsset{labeled/.style={nodes near coords, only marks, mark=*, draw=black, fill=black, thin, mark size=2, point meta=explicit symbolic}}

\usepackage{tikz}
\usetikzlibrary{
  matrix,%
  % enables the snake paths
  decorations.pathmorphing,%
  % enables bigger arrows
  decorations.markings,%
  decorations.pathreplacing,%
  decorations.text,%
  calc,%
  patterns,%
  shapes.geometric,%
  arrows,
  decorations.shapes,
  positioning,
  plotmarks,
  intersections
}
\tikzset{>=latex} % sets default arrow type in tikz
\tikzset{font=\normalsize}
\tikzset{mark size=1.5pt, mark options=thin}
\tikzset{pin distance=4pt,
  every pin edge/.style={<-, thin, shorten <= -2pt}}

\interfootnotelinepenalty=10000
\renewcommand{\thefootnote}{(\arabic{footnote})}

\usepackage{multirow, bigdelim}

% change to \usepackage[sols]{handout} to see solutions
\usepackage[sols, grading, notes]{handout}

\providecommand{\check}{}
\renewcommand{\check}{\ifmmode{\text{ \ding{51} }}\else{\ding{51} }\fi}

\newcommand\iscurrentenumi[1]{%
  \ifthenelse{\equal{#1}{\theenumi}}{}{#1}}
\setenumerate[1]{ref=\#\arabic*}
\setenumerate[2]{ref=\protect\iscurrentenumi{\theenumi}(\alph*)}
\newcommand\enumiiprefix[2]{%
  \ifthenelse{{\equal{#1}{\theenumi}}\AND{\equal{#2}{(\alph{enumii})}}}{}{}%
  \ifthenelse{{\equal{#1}{\theenumi}}\AND{\NOT{\equal{#2}{(\alph{enumii})}}}}{#2}{}%
  \ifthenelse{\equal{#1}{\theenumi}}{}{#1#2}%
}
\setenumerate[3]{ref=\protect\enumiiprefix{\theenumi}{(\alph{enumii})}\roman*}

\definecolor{purple}{rgb}{0.5,0,1.0}

\newcommand{\doedfinity}[2][\newpage]{
  \begin{problem}[#1]
    Do the Edfinity assignment ``Problem Set #2''.

    We encourage you to write down any scratch work here, as well as
    things you want your future self to remember. Nothing you write
    here will be graded; it's just for you!
  \end{problem}
}

\newcommand{\psrnewpage}{\newpage}
\newcommand{\psreflection}[2][]{

  \ifthenelse{\not\equal{#1}{nocite}}{
    \begin{enumerate}[resume]
    \item
      \begin{problem}[1in]
        Please cite any help you got on this problem set:
        \begin{itemize}[nosep]
        \item If you worked with other students, please list their names here.
        \item If you used resources other than those on our Canvas site, please list them here.
        \end{itemize}
      \end{problem}

    \end{enumerate}
  }

  \probonly{%
    #2

    \psrnewpage

    \emph{Reflection space.} It's a good habit to take a few minutes at the end of each pset to reflect on what you learned and jot down notes to your future self. Here are a few reflection questions to get you started. Nothing you write here will be graded; it's just for you!
    \begin{itemize}
    \item Take a few minutes to look back at the problems. How could you explain to someone else how to do each problem? What was your strategy, and how did you know to use that strategy? If there are problems where you don't feel able to answer this, write down which ones they are. Come back to them in a few days when we post the homework solutions!

      \vspace{1.5in}

    \item Take a look back at the learning objectives on today's worksheet; how are you feeling about each one? If there are ones you know you need more practice with, write those down here.

      \vspace{1.5in}

    \item What did you find most challenging about this material? Do you have any lingering questions about it? If so, write them down here so you don't forget them. At some point, stop by office hours or MQC to ask!

      \vfill
    \end{itemize}      
  }
}

\usepackage{fancyhdr}
\lhead{\textcolor{white}{This content is protected and may not be shared, uploaded, or distributed.}}
\rhead{}
\rfoot{\vspace*{\fill}\footnotesize\textcopyright{} \emph{Harvard Math Ma}}
\renewcommand{\headrulewidth}{0pt}
\pagestyle{fancy}
\begin{document}

\begin{center}
  \large\textsc{Problem Set 27 - Inverse Functions}\vspace{.5in}\\
  \begin{problemonly}
   \large Name: \rule{5cm}{0.15mm}\\
   \end{problemonly}
\end{center}

\begin{enumerate}
\item  \note{
    \begin{itemize}[nosep]
    \item Identify which graphs are invertible
    \item Evaluate composition of functions and inverses from table / graph
    \item Interpret $f, f', f^{-1}$ in context
    \item Match transformations of $f$ (including $f^{-1}$) with graphs
    \item Exponent rules practice: if $5^a = 3$ and $5^b = 6$, evaluate $5^{a + b}$, $5^{-a}$, $5^{2b}$, $5^{2a - b - 1}$. (This is similar to algebra we'll do next time with logs.)
    \end{itemize}
  }

  \doedfinity{27}

\item  

  Christeta is about to make roast beef in her oven. The time it will take for the beef to cook depends on the temperature of the oven; suppose $f(x)$ gives the number of minutes it will take for the beef to cook if the oven is set to $x^\circ$C. Describe in words what each of the following quantities represents; please include units, and be sure to explain what 100 represents. In addition, say whether the quantity should be positive or negative; explain how you know. 

  \begin{grading}
    2.5 points per part: 1 for the interpretation, 1 for the units, 0.5 for explaining what the 100 represents
  \end{grading}

  \begin{enumerate}
  \item
    \begin{problem}[\vfill]
      $f(100)$
    \end{problem}

    \begin{solution}
      $f(100)$ is the time (in minutes) it will take for the roast beef to cook
      if the oven is set to $100$ $^\circ$C. This will be positive,
      as it takes some amount of time to cook roast beef.
    \end{solution}

  \item
    \begin{problem}[\vfill]
      $f'(100)$
    \end{problem}

    \begin{solution}
      $f'(100)$ is the rate (in minutes per degree Celsius) 
      at which the time it takes for the roast beef to cook would change
      in response to a change in the oven temperature, when the oven temperature is
      $100$ $^\circ$C. This will be negative, as an increase in the oven temperature
      will result in a decrease in the cooking time.
      
    \end{solution}

  \item
    \begin{problem}[\vfill]
      $f^{-1}(100)$
    \end{problem}

    \begin{solution}
      $f^{-1}(100)$ is the temperature (in degrees Celsius) at which the oven
      should be set to cook the roast beef in $100$ minutes. This will be positive,
      as a negative temperature will not cook the beef.
    \end{solution}

  \end{enumerate}

\item  

  \begin{enumerate}
  \item \label{graphs-of-function-and-inverse}

    \begin{grading}
      1 point
    \end{grading}

    \begin{problem}[\vfill\newpage]
      If $f$ is an invertible function, how does the graph of $f^{-1}$
      relate to the graph of $f$? 
    \end{problem}

    \begin{solution}
      The graph of $y=f^{-1}(x)$ is a reflection across the line $y=x$ of the graph of
      $y=f(x)$.
    \end{solution}

  \item \label{based-on-gottlieb-12.1.3}

    In each part, you're given a function $f$. Sketch the graphs of $f$ and $f^{-1}$ on the same axes, using the relationship you described in \ref{graphs-of-function-and-inverse}. Please graph $f$ with a solid line and $f^{-1}$ with a dashed line so that it's clear which is which. Label 2 points on the graph of $f$ and the corresponding 2 points on the graph of $f^{-1}$.

    \begin{enumerate}
    \item

      \begin{grading}
        3 points: 0.5 for graph of $f$, 0.5 for correctly labeling 2
        points on the graph of $f$, 1 for correctly flipping it over
        $y = x$ to get the graph of $f^{-1}$, 1 for correctly flipping
        the coordinates of the points
      \end{grading}

      \begin{problem}[\vfill]
        $f(x) = 2x + 1$
      \end{problem}

      \begin{solution}
        \begin{center}
          \begin{tikzpicture}
            \begin{axis}[axis equal image, xmin=-3.75, xmax=3.75, ymin=-3.75,
              ymax=3.75, clip=false, xtick=\empty, ytick=\empty]
              \addplot+[domain=-2.375:1.375, blue]{2*x+1}
              node[above]{$f$}; %
              \addplot+[domain=-3.75:3.75, red] {(x-1)*0.5}
              node[right]{$f^{-1}$}; %
              \addplot+[domain=-3.75:3.75, dashed]{x}; %
              \addplot+[closed, blue] coordinates {(0, 1)} node[left]
              {$(0, 1)$}; %
              \addplot+[closed, blue] coordinates {(1, 3)} node[right]
              {$(1, 3)$}; %
              \addplot+[closed, red] coordinates {(1, 0)}
              node[anchor=north west] {$(0, 1)$}; %
              \addplot+[closed, red] coordinates {(3, 1)}
              node[anchor=north west] {$(3, 1)$}; %
            \end{axis}
          \end{tikzpicture}
        \end{center}
      \end{solution}

    \item

      \begin{grading}
        4 points: 1 for graph of $f$, 0.5 for correctly labeling 2
        points on the graph of $f$, 1.5 for correctly flipping it over
        $y = x$ to get the graph of $f^{-1}$, 1 for correctly flipping
        the coordinates of the points
      \end{grading}

      \begin{problem}[\vfill]
        $f$ is the function with domain $[1, \infty)$ defined by $f(x)
        = (x - 1)^2 - 2$
      \end{problem}

      \begin{solution}
        \begin{center}
          \begin{tikzpicture}
            \begin{axis}[axis equal image, clip=false, xtick=\empty, ytick=\empty]
              \addplot+[domain=1:3.5, blue]{(x-1)^2-2}
              node[above]{$f$}; %
              \addplot+[domain=-2:4.25, red] {sqrt(x+2)+1}
              node[right]{$f^{-1}$}; %
              \addplot+[domain=-5:5, dashed]{x}; %
              \addplot+[closed, blue] coordinates {(1, -2)} node[below]
              {$(1, -2)$}; %
              \addplot+[closed, blue] coordinates {(3, 2)} node[right]
              {$(3, 2)$}; %
              \addplot+[closed, red] coordinates {(-2, 1)}
              node[left] {$(-2, 1)$}; %
              \addplot+[closed, red] coordinates {(2, 3)}
              node[above] {$(2, 3)$}; %
            \end{axis}
          \end{tikzpicture}
        \end{center}
      \end{solution}

    \end{enumerate}

  \item

    \begin{grading}
      3 points, 1.5 per function
    \end{grading}

    \begin{problem}[\vfill\newpage]
      For each function in \ref{based-on-gottlieb-12.1.3}, find a formula for $f^{-1}(x)$; be sure to show your work.

      \check Make sure your formula for $f^{-1}$ is consistent with your graph!
    \end{problem}

    \begin{solution}
      
      To find the formula for $f^{-1}(x)$ algebraically, we start with
      $y = f(x)$, switch $x$ and $y$, and then solve for $y$. 
      \begin{enumerate}
      \item
        \begin{align*}
          x &= 2y + 1 \\
          x - 1 &= 2y \\
          y &= \frac{x - 1}{2} \\
          f^{-1}(x) &= \boxed{\frac{x - 1}{2}}
        \end{align*}
      \item 
        \begin{align*}
          (y-1)^2 -2 &= x \\
          (y-1)^2 &= x+2 \\
          y-1 &= \sqrt{x+2} \\
          y &= \sqrt{x+2} + 1 \\
          f^{-1}(x) &= \boxed{\sqrt{x+2} + 1}
        \end{align*}
      \end{enumerate}
    \end{solution}

  \end{enumerate}

\item 

  Let $f(x) = -\dfrac{1}{(x - 2)^2} + 5$.

  \begin{enumerate}
  \item

    \begin{grading}
      3.5 points: 1 for correct shape, 1 for correct asymptotes, 1.5 for correct description of transformations 
    \end{grading}

    \begin{problem}[\vfill]
      Sketch the graph of $f$, labeling any asymptotes. Explain how the graph of $f$ is related to the graph of $g(x) = \dfrac{1}{x^2}$. That is, if you start with the graph of $g(x) = \dfrac{1}{x^2}$, what graph transformations could you do to get the graph of $f(x)$?
    \end{problem}

    \begin{solution}
      Starting with the graph of $y=\frac{1}{x^2}$, we shift to the right $2$ units,
      reflect over the $x$-axis, then shift up $5$ units.

      \begin{center}
        \begin{tikzpicture}
          \begin{axis}[xtick=\empty, ytick=\empty, extra x ticks={2}, extra y ticks={5},
          restrict y to domain=-6:6, axis equal image]
            \addplot[domain=-3:7,samples=300] {-1/(x-2)^2 + 5};
          \end{axis}
        \end{tikzpicture}
      \end{center}
    \end{solution}

  \item

    \begin{grading}
      1 point
    \end{grading}

    \begin{problem}[\stretch{0.5}]
      Explain how you can tell $f$ isn't invertible.
    \end{problem}

    \begin{solution}
      $f$ isn't invertible because the graph doesn't pass the horizontal line test.
    \end{solution}

  \item

    \begin{grading}
      1 point
    \end{grading}

    \begin{problem}[\stretch{0.5}]
      What domain could you restrict $f$ to in order to make it invertible? Please give the largest possible domain that includes $x = 100$.
    \end{problem}

    \begin{solution}
      We could restrict $f$ to the domain $(2, \infty)$.
    \end{solution}

  \item

    \begin{grading}
      1 point
    \end{grading}

    \begin{problem}[\nstretch{0.5}]
      If $g$ is the inverse of your restricted $f$, what are the domain and range of $g$?
    \end{problem}

    \begin{solution}
      The domain of $f$ is $(2, \infty)$ and the range of $f$ is $(-\infty, 5)$. 
      So the domain of $g$ is $(-\infty, 5)$ and the range of $g$ is $(2, \infty)$.
    \end{solution}

  \end{enumerate}

\item 

  Let $f(x) = e^x$.  

  \begin{enumerate}

  \item

    \begin{grading}
      2 points: 1 for the graph, 1 for correctly labeling 2 points
    \end{grading}

    \begin{problem}[\vfill]
      Sketch the graph of $f(x)$, and label two points on the graph. 
    \end{problem}

    \begin{solution}
      \begin{center}
        \begin{tikzpicture}
          \begin{axis}[xtick=\empty, ytick=\empty]
            \addplot+[domain=-2:2]{e^x};
            \addplot[closed] coordinates {(0, 1)} node[yshift=2, left] {$(0,1)$};
            \addplot[closed] coordinates {(1, 2.718)} node[left] {$(1,e)$};
          \end{axis}
        \end{tikzpicture}
      \end{center}
    \end{solution}

  \item

    \begin{grading}
      3 points: 1 for explaining why invertible, 1 for the graph, 1 for correctly labeling 2 points
    \end{grading}

    \begin{problem}[\vfill]
      Explain how you know $f$ is invertible. Then sketch the graph of $f^{-1}$ on the axes in (a). Label two points on the graph of $f^{-1}$ that are related to the two points you labeled on the graph of $f$.
    \end{problem}

    \begin{solution}
      We know that $f$ is invertible because every $y$ value has only one $x$ value associated with it - in other words, it passes the horizontal line test. In the figure below, $y=e^x$ is given in blue, $f^{-1}$ is given in red, and $y=x$ is the dashed line.
      \begin{center}
        \begin{tikzpicture}
          \begin{axis}[xtick=\empty, ytick=\empty, axis equal image]
            \addplot+[domain=-5:2] [thick, blue]{e^x};
            \addplot[closed, blue] coordinates {(0, 1)} node[left] {$(0,1)$};
            \addplot[closed, blue] coordinates {(1, 2.718)} node[left] {$(1,e)$};
            \addplot+[domain=-5:7.4] [dashed, black]{x};
            \addplot+[domain=0:7.4] [thick, red]{ln(x)};
            \addplot[closed, red] coordinates {(1, 0)} node[right] {$(1,0)$};
            \addplot[closed, red] coordinates {(2.718,1)} node[right] {$(e,1)$};
          \end{axis}
        \end{tikzpicture}
      \end{center}
    \end{solution}

  \item

    \begin{grading}
      1 point
    \end{grading}

    \begin{problem}[\stretch{0.5}]
      What's the domain of $f^{-1}$? What is its range?
    \end{problem}

    \begin{solution}
      The domain of $f^{-1}(x)$ is \fbox{$(0, \infty)$}. Its range is \fbox{$(-\infty, \infty)$}.
    \end{solution}

  \item

    \begin{grading}
      1 point
    \end{grading}

    \begin{problem}[\stretch{0.5}]
      What's $f^{-1}(1)$?
    \end{problem}

    \begin{solution}
      Looking at our plot, we see that \fbox{$f^{-1}(1)=0$}. We can also solve this algebraically - if $e^x=1$, then $x$ must equal zero.
    \end{solution}

  \end{enumerate}

\item  
  \note{Just a reminder to prepare before we talk about $\frac{d}{dx} (b^x)$ in class.}

  \begin{grading}
    1 point
  \end{grading}

  \begin{problem}[\nstretch{0.5}]
    \textbf{Reflect Back.} Last week in class (when we talked about derivatives of exponential functions), we guessed formulas for the derivatives of $2^x$, $3^x$, and $10^x$. What were those those guesses? 
  \end{problem}

  \begin{solution}
    We guessed that $\frac{d}{dx}(2^x) \approx (0.693)2^x$, 
    $\frac{d}{dx}(3^x) \approx (1.10)3^x$, and $\frac{d}{dx}(10^x) \approx (2.30)10^x$.
  \end{solution}

\end{enumerate}
\renewcommand{\psrnewpage}{{\centering\rule{5in}{0.4pt}}}

\psreflection{For next class, read \S 13.1.}
\end{document}
